% =====================================================================
% Lab03 - Mineracao de metricas DORA
% Copia datada da Introducao e das hipoteses informais (Issue #53).
%
% Data de registro: 2026-10-07
% As hipoteses abaixo foram escritas ANTES da coleta e da analise dos
% dados (Lab03S01). Esta copia nao deve ser editada depois que os
% resultados existirem: alteracoes posteriores vao para o Overleaf e a
% Discussao compara os resultados com o texto registrado aqui.
%
% Para compilar: usar o template SBC do Overleaf
% (https://www.overleaf.com/latex/templates/sbc-conferences-template/blbxwjwzdngr),
% que fornece sbc-template.sty e sbc.bst.
% =====================================================================
\documentclass[12pt]{article}

\usepackage{sbc-template}
\usepackage{graphicx,url}
\usepackage[utf8]{inputenc}
\usepackage[T1]{fontenc}
\usepackage[brazil]{babel}

\sloppy

\title{Minerando métricas DORA em repositórios open-source:\\
o quanto podemos confiar nos \textit{proxies}?}

% TODO: confirmar a instituicao e os e-mails dos tres integrantes.
\author{Luisa Clara de Paula Lara Silva, Matheus Gaston, Mirelly Alvarenga}

\address{Curso de Engenharia de Software\\
  Laboratório de Experimentação de Software
}

\begin{document}

\maketitle

\begin{abstract}
  % Escrito na entrega final, depois dos resultados.
\end{abstract}

\begin{resumo}
  % Escrito na entrega final, depois dos resultados.
\end{resumo}

\section{Introdução}

As métricas DORA (\textit{DevOps Research and Assessment}) tornaram-se a
referência mais usada para medir o desempenho de entrega de software.
Popularizadas por \cite{forsgren2018} e revisadas anualmente pelos relatórios
\textit{Accelerate State of DevOps}, elas descrevem a entrega em duas
dimensões: \textbf{velocidade}, medida pela frequência de \textit{deploy} e pelo
\textit{lead time} das mudanças, e \textbf{estabilidade}, medida pela taxa de
falha das mudanças (\textit{change failure rate}, CFR) e pelo tempo de
recuperação após um \textit{deploy} que falhou \cite{dora2024}. Em 2024 o DORA
acrescentou uma quinta métrica, a \textit{deployment rework rate}, que mede a
proporção de \textit{deploys} não planejados feitos para corrigir problemas.
Uma das afirmações centrais desse corpo de pesquisa é que velocidade e
estabilidade não competem entre si: as equipes de melhor desempenho seriam
rápidas \emph{e} estáveis ao mesmo tempo.

Nas organizações, essas métricas costumam ser obtidas de ferramentas internas
que registram cada \textit{deploy} em produção e cada incidente. Em projetos
open-source isso não existe. O GitHub não registra ``\textit{deploys} em
produção'' nem ``falhas em produção''; ele registra \textit{releases},
\textit{tags}, \textit{commits} e execuções de \textit{workflows} de
integração contínua (CI). Quem quiser calcular métricas DORA a partir desses
dados precisa usar \textbf{\textit{proxies}}, isto é, medidas indiretas no
lugar das que não podem ser observadas. Neste trabalho, uma \textit{release}
publicada faz o papel de \textit{deploy} e uma execução de CI que falhou faz
o papel de falha.

Esses \textit{proxies} têm limitações conhecidas. Primeiro, \textbf{release
não é deploy}: uma biblioteca que publica uma versão por mês não coloca nada
em produção; quem faz isso são os projetos que a utilizam, cada um no seu
ritmo. Segundo, \textbf{CI não é produção}: uma execução de \textit{pipeline}
que falha no branch principal indica um teste quebrado, um problema de
infraestrutura ou um teste instável (\textit{flaky}), e não necessariamente um
defeito percebido pelo usuário. Terceiro, as datas disponíveis são
imperfeitas: \textit{rebase} e \textit{squash merge} reescrevem a data dos
\textit{commits} e distorcem o \textit{lead time}. Por isso, o valor de uma
métrica DORA minerada depende da definição operacional escolhida, e duas
definições igualmente razoáveis podem levar a conclusões diferentes sobre o
mesmo repositório.

O \textbf{objetivo} deste estudo é minerar as métricas DORA de repositórios
open-source populares que usam GitHub Actions, numa janela de observação de
12 meses (30/09/2025 a 30/09/2026), e avaliar o quanto esses valores são
confiáveis. Além de calcular as métricas, (i) validamos o critério automático
de \textit{release} corretiva contra uma amostra rotulada manualmente por três
avaliadores e (ii) medimos o quanto a classificação DORA de cada repositório
muda quando a definição operacional muda. O estudo é guiado pelas seguintes
questões de pesquisa (RQs):

\begin{itemize}
  \item \textbf{RQ01.} Qual a frequência de \textit{deploys} dos repositórios
    populares que usam CI/CD?
  \item \textbf{RQ02.} Qual o tempo entre um \textit{commit} e seu respectivo
    \textit{deploy}?
  \item \textbf{RQ03.} Qual a taxa de falha das mudanças entregues por esses
    repositórios?
  \item \textbf{RQ04.} Qual o tempo de recuperação após uma execução de CI/CD
    com falha?
  \item \textbf{RQ05.} Repositórios com maior frequência de \textit{deploy}
    apresentam maior ou menor taxa de falha?
  \item \textbf{RQ06.} Quais características dos repositórios estão
    associadas a um melhor desempenho DORA?
  \item \textbf{RQ07.} O quanto a classificação DORA de um repositório
    depende da definição operacional escolhida?
  \item \textbf{RQ08 (bônus).} Qual a \textit{rework rate} dos repositórios
    e em que ela difere do CFR baseado em \textit{releases}?
\end{itemize}

\subsection{Hipóteses informais}

As hipóteses abaixo foram registradas em 07/10/2026, antes da coleta e da
análise dos dados, e expressam o que o grupo \emph{espera} encontrar. Uma
cópia datada está no repositório do grupo; a Seção de Discussão as confronta
com os resultados. As categorias citadas (Elite, High, Medium, Low) seguem a
tabela de referência da disciplina.

\textbf{H1 (RQ01, frequência de \textit{deploy}).} Esperamos uma mediana
entre uma \textit{release} por mês e uma por semana (categoria Medium), com
distribuição muito assimétrica à direita. A maior parte dos projetos
populares é biblioteca ou \textit{framework} e agrupa mudanças em versões
planejadas; só uma cauda pequena, com publicação automatizada a cada
\textit{merge} (\textit{nightly builds} ou ferramentas como
\textit{semantic-release}), deve chegar a Elite.

\textbf{H2 (RQ02, \textit{lead time}).} Esperamos \textit{lead times} da ordem
de semanas, com a variante (a), por \textit{release}, sistematicamente maior
que a variante (b), por \textit{commit}, em praticamente todos os
repositórios. A variante (a) é determinada pelo \textit{commit} mais antigo
incluído na \textit{release}, e basta um \textit{commit} ``esquecido'' para
inflá-la; a (b) é dominada pela massa de \textit{commits} feitos pouco antes
da publicação. Assim, a variante (a) deve situar a maioria dos repositórios
em Medium ou Low, e a (b), em High ou Medium.

\textbf{H3 (RQ03, taxa de falha).} Esperamos CFR (a), o \textit{proxy} de CI,
com mediana entre 10\% e 25\%. As execuções consideradas são as de
\textit{push} no branch principal, e boa parte das mudanças já passou pelos
\textit{checks} do \textit{pull request} antes do \textit{merge}; ainda
assim, testes instáveis e falhas de infraestrutura mantêm a taxa acima de
zero. Para o CFR (b), o \textit{proxy} de entrega, esperamos valores de mesma
ordem, mas pouco correlacionados com o CFR (a), porque as duas variantes
medem fenômenos diferentes: quebra de \textit{pipeline} e correção rápida de
uma versão publicada.

\textbf{H4 (RQ04, tempo de recuperação).} Esperamos mediana de algumas horas
(categoria High). Em projetos ativos, uma falha no branch principal costuma
ser corrigida pelo \textit{commit} seguinte ou resolvida reexecutando o
\textit{workflow}. Esperamos também uma cauda longa e uma proporção não
desprezível de episódios censurados, vindos de \textit{workflows} secundários
(documentação, \textit{lint}, \textit{benchmarks}) que permanecem quebrados
por semanas sem bloquear o desenvolvimento.

\textbf{H5 (RQ05, velocidade $\times$ estabilidade).} Para o CFR (a),
esperamos correlação de Spearman fraca ou nula ($|\rho| < 0{,}3$), o que seria
compatível com a afirmação do DORA de que velocidade e estabilidade não são
um \textit{trade-off}. Para o CFR (b), esperamos correlação positiva, mas por
um artefato da definição: quem publica \textit{releases} com muita
frequência tem mais chance de ter qualquer \textit{release} seguida, em até 7
dias, por outra que só muda o \textit{patch}, mesmo que nenhuma delas corrija
um defeito.

\textbf{H6 (RQ06, fatores associados ao desempenho).} Esperamos que o número
de contribuidores esteja associado a maior frequência de \textit{deploy} e a
menor tempo de recuperação, porque mais pessoas produzem mais mudanças e
reagem mais rápido a quebras. Esperamos também diferenças por linguagem,
com ecossistemas de publicação automatizada (JavaScript/TypeScript, Go e
Rust) publicando com mais frequência, e repositórios mais antigos com
processos de \textit{release} mais lentos e cerimoniosos. Mesmo quando as
diferenças forem significativas após a correção de Holm, esperamos tamanhos
de efeito pequenos ($\varepsilon^2 < 0{,}06$), já que nenhum fator isolado
deve explicar muito da variação.

\textbf{H7 (RQ07, sensibilidade às definições).} Esperamos que a
classificação DORA seja sensível à definição operacional: pelo menos 30\% dos
repositórios devem mudar de categoria em algum par de combinações, com kappa
de Cohen ponderado apenas moderado (entre 0,4 e 0,6). A troca mais
instável deve ser a que usa \textit{tags} como unidade de \textit{deploy},
pois muitos projetos criam \textit{tags} sem publicar \textit{release}, o que
altera a frequência e, por consequência, o \textit{lead time}.

\textbf{H8 (RQ08, \textit{rework rate}).} Esperamos \textit{rework rate}
maior que o CFR (b), na faixa de 20\% a 35\% das \textit{releases}. A
\textit{rework rate} conta todas as \textit{releases} corretivas, inclusive as
publicadas mais de 7 dias depois da versão que corrigem e as que corrigem
outra \textit{release} corretiva; o CFR (b) só conta a \textit{release}
original quando a correção chega dentro da janela de 7 dias.

O restante do artigo está organizado da seguinte forma: a Seção~2 descreve a
metodologia (fonte de dados, funil de seleção, definições operacionais e
validação manual); a Seção~3 apresenta os resultados por RQ; a Seção~4
discute os resultados frente às hipóteses; a Seção~5 trata das ameaças à
validade; e a Seção~6 relata a replicação cruzada do estudo.

% \section{Metodologia}            % Lab03S02
% \section{Resultados}             % Lab03S03
% \section{Discussão}              % Lab03S03
% \section{Ameaças à validade}     % entrega final
% \section{Replicação cruzada}     % entrega final

\begin{thebibliography}{9}

\bibitem{forsgren2018}
Forsgren, N., Humble, J. e Kim, G. (2018).
\newblock \emph{Accelerate: The Science of Lean Software and DevOps}.
\newblock IT Revolution.

\bibitem{dora2024}
DORA (2024).
\newblock DORA's software delivery metrics: the four keys.
\newblock \url{https://dora.dev/guides/dora-metrics-four-keys/}.

\end{thebibliography}

\end{document}
